\begin{figure}[tbp]
\centering
\begin{tikzpicture}[x=1cm,y=1cm,font=\small,
  line cap=round,line join=round,>=Latex]
  \node[anchor=west,font=\small\bfseries] at (0,4.05)
    {(a) Nearly aligned blocks};
  \node at (3.50,3.48) {$f(x)=\pm x+a$};
  \fill[black!18] (1.40,2.75) rectangle (5.35,3.05);
  \fill[black!18] (1.40,1.95) rectangle (5.35,2.25);
  \draw (1.40,2.75) rectangle (5.60,3.05);
  \draw (1.15,1.95) rectangle (5.35,2.25);
  \node[anchor=east] at (1.00,2.90) {$f(B_i)$};
  \node[anchor=east] at (1.00,2.10) {$B_j$};
  \draw[densely dashed,black!55] (1.40,1.70) -- (1.40,3.18);
  \draw[densely dashed,black!55] (5.35,1.70) -- (5.35,3.18);
  \draw[<->] (1.40,1.45) -- (5.35,1.45);
  \node[align=center] at (3.38,0.90)
    {$J$: compared portion\\$|J|\ge(1-\kappa)s,\quad\kappa=o(1)$};
  \node[font=\footnotesize,align=center] at (3.38,0.20)
    {Each full block has $s+O(1)$ positions.};

  \node[anchor=west,font=\small\bfseries] at (7.00,4.05)
    {(b) Occurrences and underlying blocks};
  \node[anchor=west,font=\footnotesize] at (7.05,3.52)
    {Two disjoint occurrence matches};
  \node[draw,rounded corners=1pt,minimum width=1.02cm,
    minimum height=0.48cm] (i1) at (8.15,2.93) {$i^{(1)}$};
  \node[draw,rounded corners=1pt,minimum width=1.02cm,
    minimum height=0.48cm] (j1) at (11.65,2.93) {$j^{(1)}$};
  \node[draw,rounded corners=1pt,minimum width=1.02cm,
    minimum height=0.48cm] (i2) at (8.15,2.18) {$i^{(2)}$};
  \node[draw,rounded corners=1pt,minimum width=1.02cm,
    minimum height=0.48cm] (k1) at (11.65,2.18) {$k^{(1)}$};
  \draw (i1) -- (j1);
  \draw (i2) -- (k1);
  \draw[->] (9.90,1.75) -- (9.90,1.15)
    node[midway,right,font=\footnotesize] {forget occurrence indices};
  \node[draw,circle,inner sep=2.5pt] (j) at (8.15,0.64) {$j$};
  \node[draw,circle,inner sep=2.5pt,fill=black!12] (i) at (9.90,0.64) {$i$};
  \node[draw,circle,inner sep=2.5pt] (k) at (11.65,0.64) {$k$};
  \draw (j) -- (i) -- (k);
  \node[font=\footnotesize] at (9.90,0.02)
    {One underlying block $i$, two links};
\end{tikzpicture}
\caption{Aligned blocks and the two levels of comparison.  In (a), away
from subcomparison ends, the image of a full grid interval is displaced
by $o(s)$ from a partner grid interval; the shaded portion $J$ records a
comparison of $s-o(s)$ positions.  Reflection reverses order and inverts
letters.  The intervals and their small end discrepancies are schematic.
In (b), $i^{(r)}$ denotes the $r$th occurrence of block $i$.  Occurrence
joins form a partial matching, but forgetting the occurrence indices
can give several links at one underlying block.}
\label{fig:aligned-blocks}
\end{figure}
